\section{Setta del Sangue}\label{setta-del-sangue}

Tags: Organizzazione Creatore: Lorenzo

\section{Setta del Sangue}\label{setta-del-sangue-1}

\begin{center}\rule{0.5\linewidth}{0.5pt}\end{center}

\begin{center}\rule{0.5\linewidth}{0.5pt}\end{center}

\begin{figure}
\centering
\includegraphics{simbolo.png}
\caption{simbolo.png}
\end{figure}

Informazioni Generali

Tipo: Setta

Struttura: Organizzazione segreta

Regione: Valtara

Fondatore: Sconosciuto

Membri: Sconosciuto

Alleati: Sconosciuto

Nemesi: Ordine dei Paladini di San Francesco; Gilda del Protettori

\subsection{1. Descrizione Generale}\label{descrizione-generale}

\begin{center}\rule{0.5\linewidth}{0.5pt}\end{center}

\begin{figure}
\centering
\includegraphics{generate-an-image-depicting-the-exterior-of-il-pane-quotidiano-newspaper-headquarters-in-the-medie.png}
\caption{generate-an-image-depicting-the-exterior-of-il-pane-quotidiano-newspaper-headquarters-in-the-medie.png}
\end{figure}

\emph{Il Pane Quotidiano} è un rinomato giornale nel mondo di Valtara,
noto per la sua dedizione nell'offrire una copertura completa e accurata
degli avvenimenti in . Fondata con l'obiettivo di soddisfare la sete di
conoscenza e la curiosità dei suoi lettori, questa organizzazione
giornalistica ha stabilito la sua sede principale nella maestosa città
di Kos.

\begin{quote}
{[}Quote by or about organization, possibly from founding member{]}
\end{quote}

\subsection{2. Storia}\label{storia}

\begin{center}\rule{0.5\linewidth}{0.5pt}\end{center}

La storia della Setta del Sangue è per lo più sconosciuta e oscura. In
tempi antichi, la setta era originariamente una congregazione di
individui alla ricerca di potere e dominio attraverso lo spionaggio e la
corruzione. Tuttavia, la loro vita prese una piega sinistra quando si
imbatterono nelle leggende del demone Naskirophis.

Queste antiche profezie riguardanti il demone suscitarono l'interesse
della setta, che decise di dedicarsi al culto di Naskirophis. Per
generazioni, questi devoti hanno trascorso le loro esistenze alla
ricerca di metodi per evocare il demone e ottenere il suo potere,
credendo che avrebbe portato loro dominio e controllo sul mondo.

La Setta del Sangue divenne maestra nell'arte dell'inganno e
dell'infiltrazione. Durante la Guerra del Sangue, un conflitto tra
l'Ordine dei Paladini di San Francesco e forze oscure che cercavano di
liberare Naskirophis, la Setta del Sangue si infiltrò segretamente tra i
ranghi dei Paladini, sfruttando i loro poteri oscuri per soccombere la
lotta dall'interno.

La guerra si concluse con la vittoria dei Paladini, che riuscirono a
imprigionare Naskirophis in una prigione dimensionale, mettendo fine al
conflitto noto come la Guerra del Sangue. Nonostante la sconfitta, la
Setta del Sangue sopravvisse, mantenendo viva la fiamma del culto del
demone.

Negli ultimi 20 anni, la Setta del Sangue è rimasta nell'ombra,
conservando la sua devozione al demone Naskirophis e cercando nuove
opportunità per diffondere il suo culto di oscurità e corruzione.
Recentemente, il demone è stato inaspettatamente liberato dalla sua
prigione dimensionale, scatenando un nuovo periodo di terrore. La Setta
del Sangue ha cercato di riportare agli antichi albori il culto del
demone, ma è stata fermata dai Protettori, una gilda di avventurieri
devoti alla giustizia.

Sebbene il demone sia stato sconfitto, la Setta del Sangue rimane una
minaccia costante, sempre pronta a cercare nuove opportunità per
perseguire il suo obiettivo di creare uno stato basato sul dominio
demoniaco.

\subsection{3. Valori}\label{valori}

\begin{center}\rule{0.5\linewidth}{0.5pt}\end{center}

\begin{enumerate}
\def\labelenumi{\arabic{enumi}.}
\tightlist
\item
  \textbf{Devozione a Naskirophis:} La Setta considera Naskirophis come
  una divinità suprema, dedicando le loro vite a servirlo e ottenere il
  suo favore. La devozione a questa entità demoniaca è al centro dei
  loro valori.
\item
  \textbf{Sacrifici di Sangue:} La Setta pratica sacrifici rituali di
  sangue come parte integrante della loro adorazione. Credono che questi
  sacrifici (spesso di infanti) siano un mezzo per guadagnare potere o
  favori dal loro signore oscuro.
\item
  \textbf{Conversione:} La Setta cerca attivamente di convertire altre
  persone al loro credo, diffondendo il culto di Naskirophis. Vedono la
  conversione come un passo verso l'espansione del dominio di
  Naskirophis nel mondo.
\item
  \textbf{Oscurità e Caos:} I membri della Setta abbracciano l'oscurità
  e il caos come elementi essenziali della loro fede. Credono che solo
  attraverso il caos e la destabilizzazione possano raggiungere i loro
  obiettivi.
\item
  \textbf{Potere Personale:} Molti membri della Setta sono motivati
  dalla ricerca del potere personale e dell'immortalità, credendo che
  l'adorazione di Naskirophis sia il mezzo per raggiungere questi scopi.
\end{enumerate}

\subsection{4. Cultura}\label{cultura}

\begin{center}\rule{0.5\linewidth}{0.5pt}\end{center}

La cultura della Setta del Sangue è intrisa di oscurità e fanatismo. I
suoi membri venerano il demone Naskirophis come una divinità,
sacrificando innocenti neonati per ottenere potere. La lealtà alla setta
è considerata più importante di qualsiasi altra cosa, e il segreto e la
clandestinità sono fondamentali. I membri sono indottrinati fin dalla
giovane età, e la violenza è un mezzo comune per raggiungere i loro
scopi. La setta mira a creare uno stato demoniaco e caotico, dove il
male regna sovrano. La sua cultura è basata su ideali malvagi e devoti
al demone, ed è profondamente avversa alla pace e alla giustizia.

\subsubsection{4.1 Educazione \&
Apprendistato}\label{educazione-apprendistato}

\begin{center}\rule{0.5\linewidth}{0.5pt}\end{center}

Diventare membri della Setta del Sangue è un processo che richiede
dedizione assoluta e profonda devotizione al demone Naskirophis. Non si
tratta di un'organizzazione aperta o accessibile a chiunque, ma è
riservata a coloro che dimostrano una fede fanatica nell'entità
demoniaca e sono disposti a compiere atti malvagi per servirla.

In genere, il reclutamento nella Setta del Sangue avviene attraverso un
processo di selezione segreto e rigoroso. Gli individui, spesso
reclutati fin ragazzi, devono dimostrare la loro lealtà e devotamente
attraverso azioni e sacrifici che favoriscono gli interessi della setta.
Questi possono includere atti di corruzione, violenza e persino omicidi.
L'obiettivo principale è dimostrare la fedeltà assoluta a Naskirophis.

Una volta che un individuo ha dimostrato la sua devozione, può essere
invitato a partecipare ai rituali segreti della setta e ad apprendere le
arti della politica e della dialettica, in modo da potersi infiltrare
negli ambienti altolocati della società Nord-Valtarese. Tuttavia, è
importante sottolineare che la Setta del Sangue è un'organizzazione
clandestina e illegale, e coloro che cercano di entrarvi sono
considerati criminali dalla società.

\subsection{5. Attività}\label{attivituxe0}

\begin{center}\rule{0.5\linewidth}{0.5pt}\end{center}

La Setta del Sangue, devota a Naskirophis, ha come primario obiettivo
l'adorazione e l'ottenimento del favore del loro oscuro signore. Il loro
culto ruota attorno a sacrifici rituali di sangue, con l'intento di
ricevere poteri sovrumani o doni oscuri in cambio della loro lealtà.
Tuttavia, il loro zelo non si ferma qui: cercano anche di diffondere il
culto di Naskirophis tra le popolazioni locali, convertendo le anime
smarrite al loro credo. L'obiettivo finale è consolidare il potere di
Naskirophis nel mondo, influenzando decisioni politiche, seminando
terrore e caos, e portando la volontà oscura dell'entità demoniaca nella
società in cui operano.
